\documentclass[letterpaper,10pt]{article}

\usepackage[hyphens]{url}
\usepackage{graphicx}
\usepackage{amsmath}
\usepackage{amssymb}
\newcommand{\toprule}{\hline}
\newcommand{\midrule}{\hline}
\newcommand{\bottomrule}{\hline}

\graphicspath{{../runs/kd_failure_analysis/figures/}}
\urlstyle{rm}
\def\UrlFont{\rm}
\frenchspacing
\setlength{\pdfpagewidth}{8.5in}
\setlength{\pdfpageheight}{11in}

\title{\textbf{CoughKD-ShiftAudit: Failure Cartography for Ultra-Light Cough Audio Distillation under Dataset Shift}}
\author{Anonymous Authors\\Anonymous Institution}
\date{}

\begin{document}
\maketitle

\begin{abstract}
Large audio models can provide useful representations for cough and respiratory sounds, but deploying them on edge devices requires compact students. Knowledge distillation (KD) is a natural compression tool, yet its reliability under cough dataset shift remains unclear. We introduce \textsc{CoughKD-ShiftAudit}, a deployment-oriented audit protocol for ultra-light cough audio distillation. Using Coswara as the source dataset and COUGHVID plus Tos COVID-19 as independent public external targets, we evaluate CE-only training, vanilla KD, target-consistency distillation, shortcut-suppressed KD, disagreement-gated KD, and probe-adversarial training across multiple seeds. Beyond aggregate AUROC, we report calibration, task and domain probes, high-specificity sensitivity, model size, latency, metadata-defined target slices, bootstrap target-subset stability, overlap control, and clip-vs-subject stress tests. Current evidence shows that the best aggregate KD variant improves COUGHVID external macro AUROC by only $+0.0007$ over source-only continuation; on Tos COVID-19, CE-only training is the post-hoc best method and vanilla KD is negative. Small Virufy stress targets further change method rankings across preprocessing and evaluation units. These findings suggest that cough audio distillation under dataset shift requires failure-oriented auditing before deployment, and that shortcut or calibration improvements alone should not be interpreted as external robustness.
\end{abstract}

\section{Introduction}

Cough is a frequent symptom of respiratory disease and can be recorded with commodity microphones at negligible marginal cost. Large crowdsourced corpora such as COUGHVID and Coswara have made cough audio a realistic research substrate for COVID-19 and respiratory-condition screening \cite{orlandic2021coughvid,bhattacharya2023coswara}. However, the practical value of a cough classifier depends on more than in-domain accuracy: it must tolerate uncontrolled microphones, background speech, label noise, demographic shift, and low-resource deployment \cite{han2022sounds}.

Recent audio foundation models, including AST, PANNs, PaSST, HTS-AT, AudioMAE, and BEATs, provide strong transferable representations for audio classification \cite{gong2021ast,kong2020panns,koutini2022passt,chen2022htsat,baade2022maeast,chen2023beats}. Their capacity is useful for noisy cough data, but their memory and latency profiles are poorly matched to mobile screening or wearable health scenarios. This motivates a focused question:

\begin{quote}
\emph{Can a compact cough screening network retain the cross-dataset performance and calibration of an audio foundation-model teacher while achieving an order-of-magnitude reduction in computation?}
\end{quote}

We propose \textsc{CoughKD}, a compression and distillation framework for cough audio classification. The teacher is a high-capacity audio foundation model fine-tuned with conservative data protocols; the student is a compact model selected from deployment-friendly architectures such as MobileNetV3, EfficientNet-B0, BC-ResNet, and ECAPA-small. Distillation is performed at four complementary levels: softened responses, intermediate representations, attention maps, and sample relations. The design is intentionally empirical and auditable: each component is paired with an ablation and a failure mode.

The intended contributions are:

\begin{itemize}
    \item A reproducible teacher--student protocol for cough audio screening that separates in-domain performance from cross-dataset generalization.
    \item A multi-level distillation objective combining response, feature, attention, and relation transfer for compact respiratory audio models.
    \item A benchmark plan covering more than twenty teacher, student, classical, and distilled baselines under identical subject-disjoint splits.
    \item A deployment-oriented evaluation suite reporting AUROC/AUPRC/F1 together with calibration error, latency, parameter count, FLOPs, and model size.
\end{itemize}

The current document is a paper scaffold and research protocol. It uses target values only to define acceptance criteria for the study. No target value should be presented as an experimental result until reproduced under the protocol in Section~\ref{sec:experiments}.

\section{Related Work}

\paragraph{Cough and respiratory audio datasets.}
COUGHVID provides a large-scale crowdsourced cough corpus with expert annotation for a subset of recordings \cite{orlandic2021coughvid}. Coswara extends cough-only screening with breathing, phonation, speech, demographics, symptoms, and SARS-CoV-2 status \cite{bhattacharya2023coswara}. Cambridge COVID-19 Sounds and the DiCOVA challenges highlighted a central methodological issue: cough models can perform strongly under narrow splits while degrading under realistic external validation \cite{han2022sounds}. ICBHI 2017 is not a cough dataset, but it remains useful for respiratory-sound transfer and robustness analysis.

\paragraph{Cough-based disease screening.}
Earlier cough systems used MFCC, spectral descriptors, bag-of-words, SVMs, random forests, and shallow neural networks \cite{pavel2023bow,pahar2021cough}. Later systems adopted CNNs, LSTMs, spectrogram image classifiers, and ensemble learning \cite{hussain2024cough2covid}. Across this literature, AUROC, sensitivity, specificity, F1, and subgroup performance are more informative than accuracy alone because cough datasets are often imbalanced and clinically asymmetric. Recent high reported accuracies should be interpreted carefully unless subject-level splits, external datasets, and threshold selection policies are explicit.

\paragraph{Audio foundation models.}
PANNs demonstrated that CNNs pretrained on AudioSet transfer well across audio tasks \cite{kong2020panns}. AST adapted the Vision Transformer to spectrogram patches and showed the strength of attention-only audio classification \cite{gong2021ast}. PaSST improved transformer training efficiency through Patchout \cite{koutini2022passt}. HTS-AT introduced hierarchical token-semantic structure for sound classification and detection \cite{chen2022htsat}. BEATs used acoustic tokenizers for iterative self-supervised audio pretraining and achieved strong AudioSet and downstream performance \cite{chen2023beats}. These models are strong teachers for cough screening, but they are not ideal endpoint models for low-power devices.

\paragraph{Efficient audio students.}
MobileNetV3 and EfficientNet are established compact vision backbones that transfer naturally to spectrogram inputs \cite{howard2019mobilenetv3,tan2019efficientnet}. BC-ResNet was designed for efficient keyword spotting and is attractive for tiny audio classifiers \cite{kim2021bcresnet}. ECAPA-TDNN is stronger for speaker embeddings than disease labels, but its channel attention and temporal pooling make it a useful compact baseline for cough representations. The main empirical question is not whether small models run faster, but whether they retain teacher robustness under dataset shift.

\paragraph{Knowledge distillation.}
Response distillation transfers softened teacher logits \cite{hinton2015distilling}. Feature distillation aligns hidden representations \cite{romero2015fitnets}. Attention transfer encourages similar saliency or activation distributions \cite{zagoruyko2017attention}. Relation and contrastive distillation preserve pairwise or class-structured geometry \cite{tian2020crd,khosla2020supcon}. Recent acoustic monitoring work also shows KD can support efficient audio detection pipelines outside standard vision/NLP settings \cite{priebe2024speechkd}. \textsc{CoughKD} combines these ideas because cough labels are noisy and sparse: response KD captures class similarity, feature KD transfers acoustic abstractions, attention KD regularizes time-frequency focus, and relation KD stabilizes representation geometry across domains.

\section{CoughKD-ShiftAudit}
\label{sec:method}

\subsection{Problem Setup}

Let $\mathcal{D}_s=\{(x_i,y_i,u_i)\}$ be a source cough dataset with recording $x_i$, label $y_i$, and subject identifier $u_i$. Let $\mathcal{D}_t$ be a target dataset used only for external evaluation and stress testing. The goal is not to maximize a single in-domain score, but to audit whether a compact student $f_s$ distilled from a high-capacity teacher $f_t$ remains reliable under $\mathcal{D}_s \rightarrow \mathcal{D}_t$ shift.

We treat the source split as subject-disjoint. Target labels are used for evaluation only, not for target-supervised model selection. When a target contains multiple clips from one subject or original recording, we report both clip-level and subject-level views when possible.

\subsection{Teacher and Student}

The current teacher is PANNs CNN14 at 16 kHz, initialized from AudioSet-pretrained weights and adapted to the source label space \cite{kong2020panns}. The student is a compact depthwise CNN with 20.7K parameters. This gives an extreme compression ratio of about 3912x relative to the 81M-parameter teacher. The teacher/student pair is intentionally simple: the paper audits whether even standard KD assumptions are reliable before expanding to a larger teacher-by-student grid.

\subsection{Audited Training Variants}

We compare a bounded set of training variants:
\begin{itemize}
    \item CE-only student training;
    \item vanilla response KD;
    \item source-only continuation;
    \item target-consistency KD;
    \item confidence-gated target-consistency KD;
    \item shortcut-suppressed KD;
    \item disagreement-gated KD;
    \item probe-adversarial KD.
\end{itemize}

These variants are not presented as a new SOTA architecture search. They form an audit set for common KD assumptions: that teacher probabilities help, that target consistency helps, that shortcut suppression helps, and that domain readability reduction helps.

\subsection{Audit Evidence}

\textsc{CoughKD-ShiftAudit} reports:
\begin{itemize}
    \item external macro AUROC, COVID AUROC, and macro AUPRC;
    \item ECE, Brier score, and sensitivity at 95\% specificity;
    \item domain probe AUC and task probe AUC over student representations;
    \item parameter count, checkpoint size, and latency;
    \item metadata-defined COUGHVID slices;
    \item bootstrap target-subset stability;
    \item target-unlabeled guard stress tests;
    \item clip-vs-subject aggregation on segmented targets;
    \item source-target overlap control for external and auxiliary manifests.
\end{itemize}

\subsection{Overlap Control}

External validation claims are invalid if the target overlaps with the source. Before reporting any new external or auxiliary target, we run a manifest overlap audit against Coswara. The audit checks exact recording IDs, subject IDs, paths, path stems, path tails, shared identifier-like tokens, and optionally SHA-256 audio hashes. DiCOVA is treated as a Coswara-related auxiliary protocol unless overlap is explicitly controlled.

\section{Experiments}
\label{sec:experiments}

\paragraph{Source and targets.}
Coswara cough-heavy and cough-shallow recordings form the source training set. COUGHVID is the largest public external target in our current audit. Tos COVID-19 is added as a second independent public external target; its 2021 official test split provides RT-PCR positive/negative labels and OGG cough clips collected by mobile phone \cite{buenosaires2021toscovid}. Virufy original and segmented clips are included only as tiny public stress targets, not as strong external validation cohorts. UK COVID-19 Vocal Audio and Cambridge COVID-19 Sounds / ComParE CCS remain high-priority larger follow-up targets, while DiCOVA is reserved as a Coswara-derived auxiliary protocol with overlap control.

\paragraph{Evaluation metrics.}
We report macro one-vs-rest AUROC, COVID one-vs-rest AUROC, macro AUPRC, ECE, Brier score, sensitivity at 95\% specificity, domain probe AUC, task probe AUC, checkpoint size, latency, and method-selection stability. Each primary COUGHVID result is averaged across seeds 7, 11, and 23.

\paragraph{Stress tests.}
COUGHVID is further analyzed with metadata-defined slices and bootstrap target subsets. Tos COVID-19 tests whether conclusions persist on a second independent public COVID cough target. Virufy segmented is evaluated both at the clip level and after subject-level probability aggregation. These stress tests ask whether the same KD method remains best when the target composition or evaluation unit changes.

\paragraph{Guard evaluation.}
We evaluate whether target-unlabeled signals can select the best KD strategy. Candidate signals include target confidence, entropy, predicted label distribution, and source-relative prediction shift. A guard is considered useful only if it selects a method that improves external metrics without increasing probe or calibration risk.

\section{Results}
\label{sec:results}

Table~\ref{tab:coughvid_external_audit} summarizes the main COUGHVID external audit. The best aggregate method, confidence-gated target-consistency KD, improves external macro AUROC by only $+0.0007$ over source-only continuation. Vanilla KD improves COVID AUROC but reduces macro AUROC. Shortcut-suppressed and disagreement-gated variants reduce some shortcut or calibration indicators but do not improve external macro AUROC. This is the central negative result: the current KD variants do not support a robust method-superiority claim.

\input{../runs/kd_failure_analysis/paper_ready_tables/paper_ready_tables.tex}

Figure~\ref{fig:protocol} shows the audit protocol. Figures~\ref{fig:macro_ece}--\ref{fig:bootstrap} illustrate metric disagreement and target-subset sensitivity.

\begin{figure}[t]
    \centering
    \includegraphics[width=\linewidth]{shift_audit_protocol.pdf}
    \caption{CoughKD-ShiftAudit protocol. Large-teacher to ultra-light-student transfer is audited with external metrics, calibration, probes, slices, bootstrap subsets, overlap checks, and efficiency evidence.}
    \label{fig:protocol}
\end{figure}

\begin{figure}[t]
    \centering
    \includegraphics[width=0.48\linewidth]{external_macro_vs_ece.pdf}
    \includegraphics[width=0.48\linewidth]{external_macro_vs_domain_probe.pdf}
    \caption{Metric disagreement on COUGHVID. Better calibration or lower domain readability does not necessarily imply better external macro AUROC.}
    \label{fig:macro_ece}
\end{figure}

\begin{figure}[t]
    \centering
    \includegraphics[width=0.48\linewidth]{slice_delta_gap.pdf}
    \includegraphics[width=0.48\linewidth]{bootstrap_delta_distribution.pdf}
    \caption{Target-composition stress tests. Slice-level and bootstrap-subset behavior reveals larger local variation than aggregate COUGHVID AUROC suggests.}
    \label{fig:bootstrap}
\end{figure}

Table~\ref{tab:target_unit_stress} further shows that target choice and evaluation unit change conclusions. COUGHVID selects confidence-gated TCD with a tiny gain, but Tos COVID-19 selects CE-only as the post-hoc best method and assigns negative deltas to both vanilla KD and the guard-selected TCD variant. Virufy original selects CE-only; Virufy segmented clip-level selects vanilla KD; Virufy segmented subject-level selects probe-adversarial KD. The guard does not reliably track these post-hoc best methods. Because Virufy has only 16 effective source subjects, it remains a stress target rather than strong external validation. The two main external targets, COUGHVID and Tos, support the narrower conclusion that KD behavior is weak and target-sensitive rather than robustly superior.

\section{Discussion and Review Notes}

\paragraph{Why distillation is the right main story.}
Patient voiceprint recognition from cough is scientifically interesting but depends heavily on identity labels and data-use permissions. In contrast, teacher-to-student compression can be evaluated on existing cough datasets with standard classification protocols. It creates a cleaner top-conference story: strong teacher, compact deployable student, rigorous generalization and efficiency evidence.

\paragraph{Main threat to validity.}
The biggest risk is dataset confounding. COVID labels may correlate with microphone type, geography, recording date, language, or symptom-reporting behavior. Distillation can transfer these shortcuts from teacher to student. The paper must therefore emphasize external validation, subgroup reporting, and calibration rather than only leaderboard-style accuracy.

\paragraph{Expected reviewer concerns.}
Reviewers will likely ask whether the method is more than combining known KD losses. The response must be empirical and domain-specific: cough audio has label noise, segment uncertainty, and long-recording aggregation; the contribution is the validated distillation recipe under realistic respiratory-audio constraints. If full KD does not beat simpler KD, the paper should pivot to a careful negative result or a narrower claim.

\paragraph{Clinical caution.}
The system should be described as a screening research model, not a diagnostic device. Any medical claim requires prospective clinical validation and regulatory review.

\section{Conclusion}

We presented \textsc{CoughKD-ShiftAudit}, a failure-oriented audit protocol for ultra-light cough audio distillation under dataset shift. Current experiments show that standard and shortcut-aware KD variants provide weak, unstable, or negative external gains across COUGHVID and Tos COVID-19, while small Virufy stress targets change method rankings across preprocessing and evaluation units. These findings caution against single-metric deployment claims for cough KD. The next step is to strengthen the audit with larger external targets such as UK COVID-19 Vocal Audio or Cambridge COVID-19 Sounds / ComParE CCS, or to extend the protocol to a broader cough transfer setting.


\bibliography{references}
\bibliographystyle{plain}

\end{document}
